%----------------------------------------------------------------------------------------
%	                            PORTADA
%----------------------------------------------------------------------------------------
% \makeportada{}{Daniel Carbajales Soto}{ 1° CFGS Automatización y Robótica Industrial}{}

%----------------------------------------------------------------------------------------
%                                    HEADER 
%----------------------------------------------------------------------------------------
% \header{}{Daniel Carbajales Soto}{1° CFGSARI}
%----------------------------------------------------------------------------------------
%                               TABLE OF CONTENTS
%----------------------------------------------------------------------------------------
% \maketoc
%----------------------------------------------------------------------------------------
%	                            TABLES
%----------------------------------------------------------------------------------------
% this is just a demo in case i need to make a table
%
%
%
% \begin{tabularx}{\linewidth}{|C{2cm}|Y|Y|C{2cm}|}
% \hline
% \makecell{Header 1} & \makecell{Header 2 \\ with linebreak} & \makecell{Header 3} & \makecell{Header 4} \\
% \hline
% Row 1 & This is a long cell text that automatically wraps in the Y column & Another long cell & Short \\
% \hline
% Row 2 & \multirow{2}{*}{Multirow text} & More text & Short \\
% \cline{1-1} \cline{3-4}
% Row 3 & & Extra text & Short \\
% \hline
% \end{tabularx}


%----------------------------------------------------------------------------------------
%	                            DOCUMENT
%----------------------------------------------------------------------------------------

\addcontentsline{toc}{section}{\textcolor{waveBlue2}{Ejercicio 8}}
\section*{\textcolor{waveBlue2}{8. }Completa el siguiente texto con las palabras que faltan.}

La \underline{\textcolor{springViolet1}{tecnologia de la información}} (IT) y la \underline{\textcolor{springViolet1}{tecnologia de operación}} (OT) son áreas distintas dentro de una organización; 
la IT se enfoca en la gestión de la información y sistemas empresariales, y la OT, en el control de procesos físicos operativos en sectores
como \underline{\textcolor{springViolet1}{industria}} y \underline{\textcolor{springViolet1}{energía}}. Estas áreas están convergiendo debido a la 
digitalización, el auge del \underline{\textcolor{springViolet1}{IoT}} y la necesidad de integrar los datos. Dicha convergencia busca unir 
tecnologías y procesos para aumentar la eficiencia operativa, la \underline{\textcolor{springViolet1}{agilidad}} y una toma de decisiones más informada,
con una gestión centralizada de la seguridad y los datos. Para enfrentar estos desafíos, las empresas están formando \underline{\textcolor{springViolet1}{equipos multidisciplinares}}
y asociaciones estratégicas.



%%%%%%%%%%%%%%%%%%%%%%%%%%%%%%%%%%%%%%%%%%%%%%%%%%%%%%%%%%
%%%%%%%%%%%%%%%%%%%% bibliography %%%%%%%%%%%%%%%%%%%%%%%%
% \nocite{*}                                 %%%%%%%%%%%%%
% \printbibliography                         %%%%%%%%%%%%%
%%%%%%%%%%%%%%%%%%%%%%%%%%%%%%%%%%%%%%%%%%%%%%%%%%%%%%%%%%
%%%%%% Last page in case there is no bibliography %%%%%%%%
% \label{Last Page}                           %%%%%%%%%%%%%
%%%%%%%%%%%%%%%%%%%%%%%%%%%%%%%%%%%%%%%%%%%%%%%%%%%%%%%%%%

